% TOP_PROBLEM: {"problemId": "problem.equivalence-of-worst-case-and-average-case-hardness-for-np-excluding-heuristica", "canonicalTitle": "Equivalence of Worst-Case and Average-Case Hardness for NP (Excluding Heuristica)", "releaseRank": 52, "primaryDomain": "theoretical_computer_science", "primaryDomainLabel": "Theoretical computer science", "displayStatus": "open", "releaseStatus": "open"}
% !TeX program = lualatex
% Definitions and sources retained from research_batch_51_54.tex; research is in GPT_6_Astra_Ultra.
\documentclass[11pt]{article}
\usepackage[margin=25mm]{geometry}
\usepackage{fontspec}
\setmainfont{Latin Modern Roman}
\usepackage{amsmath,amssymb,mathtools}
\usepackage{unicode-math}
\setmathfont{Latin Modern Math}
\usepackage{enumitem,needspace}
\usepackage{xurl,hyperref}
\hypersetup{colorlinks=true,urlcolor=blue}
\setcounter{secnumdepth}{1}
\setlength{\parindent}{0pt}
\setlength{\parskip}{5pt}
\setlength{\emergencystretch}{3em}
\setlist[itemize]{leftmargin=3em,itemsep=5pt,topsep=4pt,parsep=0pt}
\allowdisplaybreaks
\newcommand{\N}{\mathbb{N}}
\newcommand{\Z}{\mathbb{Z}}
\newcommand{\Q}{\mathbb{Q}}
\newcommand{\R}{\mathbb{R}}
\newcommand{\C}{\mathbb{C}}
\newcommand{\F}{\mathbb{F}}
\newcommand{\A}{\mathbb{A}}
\newcommand{\PP}{\mathbb P}
\newcommand{\E}{\mathbb{E}}
\newcommand{\eps}{\varepsilon}
\newcommand{\dd}{\,\mathrm d}
\newcommand{\sref}[2]{\hyperlink{source:#1:#2}{[S#2]}}
\newcommand{\eref}[2]{\hyperlink{extra:#1:#2}{[E#2]}}
% Turn the next switch off for a shorter reading copy. The quotations preserve
% every exactTarget field, including qualifications omitted from a short summary.
\newif\ifshowcatalogue
\showcataloguetrue
\newcommand{\cataloguescope}[1]{\ifshowcatalogue
 \paragraph{Exact scope in the supplied catalogue.}
 \begin{quote}\small #1\end{quote}\fi}
\begin{document}
\setcounter{section}{51}
% ================================================================
% problemId: problem.equivalence-of-worst-case-and-average-case-hardness-for-np-excluding-heuristica
% source releaseRank: 52; source releaseStatus: open
% exactTarget SHA-256: 163c4bd826d5eb30d90067247a03b4f8037adcce0ac03cc5d63964f79fec3e18
\clearpage
\section{\texorpdfstring{Equivalence of Worst-Case and Average-Case Hardness for NP (Excluding Heuristica)}{Equivalence of Worst-Case and Average-Case Hardness for NP (Excluding Heuristica)}}\label{52}
\subsection{Definitions and mathematical statement}
A distributional problem is a pair $(L,D)$, where $L$ is a language and $D=(D_n)$ is an efficiently described input distribution ensemble. In Levin's average-time convention, a correct deterministic decider $A$ is average-polynomial if for some $\varepsilon>0$,
\[
 \mathbb E_{x\sim D_n} T_A(x)^\varepsilon\le\operatorname{poly}(n).
\]
Write $\mathsf{AvgP}$ for this class. Distributional NP uses $L\in\mathsf{NP}$ and a prescribed efficient distribution class. The target is $\mathsf{NP}\nsubseteq\mathsf{P/poly}\Rightarrow\mathsf{DistNP}\nsubseteq\mathsf{AvgP}$. The catalogue leaves polynomial-time computable versus samplable ensembles implicit; these are distinct definitions and must not be interchanged without the relevant reduction theorem. The added Bogdanov--Trevisan survey gives both conventions and distinguishes errorless algorithms from heuristics allowed to return incorrect answers.
\par\smallskip\textit{Formulation and concept references:} \sref{52}{1}, \eref{52}{1}.
\cataloguescope{Prove or disprove that if NP is hard in the worst case (NP not in P/poly), then NP is hard on average (DistNP not in AvgP).}
\subsection{Short English statement}
Does strong worst-case hardness in NP necessarily create an efficiently distributed NP problem that is hard on average?
\subsection{Sources}
\begin{itemize}\raggedright
\item[\textnormal{[S1]}]\hypertarget{source:52:1}{} A. Bogdanov, L. Trevisan, SIAM J. Comput. 36(4):1119-1159, 2006.
\url{https://doi.org/10.1137/050634629}.
{\small\textit{Catalogue formulation source.}}
\item[\textnormal{[S2]}]\hypertarget{source:52:2}{} Cryptography meets worst-case complexity: Optimal security and more from iO and worst-case assumptions.
\url{https://eccc.weizmann.ac.il/report/2025/072/download}.
{\small\textit{Catalogue research/status reference; not a proof endorsement.}}
\item[\textnormal{[S3]}]\hypertarget{source:52:3}{} Hardness of Computing Nondeterministic Kolmogorov Complexity.
\url{https://eccc.weizmann.ac.il/report/2025/203/}.
{\small\textit{Catalogue research/status reference; not a proof endorsement.}}
\item[\textnormal{[E1]}]\hypertarget{extra:52:1}{} Andrej Bogdanov and Luca Trevisan, Average-Case Complexity, Foundations and Trends in Theoretical Computer Science 2 (2006), 1--106.
\url{https://andrejb.net/pubs/average-now.pdf}.
{\small\textit{Added primary source: Definitions of distributional problems and average-polynomial time. Consulted 24 September 2026.}}
\item[\textnormal{[E2]}]\hypertarget{extra:52:2}{} Andrej Bogdanov and Luca Trevisan, \emph{On Worst-Case to Average-Case Reductions for NP Problems}, SIAM Journal on Computing 36(4) (2006), 1119--1159, DOI 10.1137/S0097539705446974. See also the precise nonadaptive formulation in Theorem 7.5 of [E1].
\url{https://doi.org/10.1137/S0097539705446974}.
{\small\textit{Added research source: Corrected identifier for the paper intended by [S1]; nonadaptive reduction barrier. The original [S1] is preserved. Consulted 24 September 2026.}}
\item[\textnormal{[E3]}]\hypertarget{extra:52:3}{} Rahul Ilango and Alex Lombardi, \emph{Cryptography meets worst-case complexity: Optimal security and more from iO and worst-case assumptions}, ECCC TR25-072 (2025), abstract and stated cryptographic hypotheses.
\url{https://eccc.weizmann.ac.il/report/2025/072/}.
{\small\textit{Added research source: The additional obfuscation and fine-grained assumptions are retained, not treated as consequences of the target premise. Consulted 24 September 2026.}}
\item[\textnormal{[E4]}]\hypertarget{extra:52:4}{} Jinqiao Hu, Zhenjian Lu and Igor Oliveira, \emph{Hardness of Computing Nondeterministic Kolmogorov Complexity}, ECCC TR25-203, revision 2, 5 July 2026, Section 1.2.3 and Appendix B, including Corollary 59.
\url{https://eccc.weizmann.ac.il/report/2025/203/revision/2/download}.
{\small\textit{Added research source: Recent meta-complexity implications and the remaining gap-version obstacle; not identified with the deterministic average-time assertion. Consulted 24 September 2026.}}
% END RESEARCH ADDITION REFERENCES-52
\end{itemize}

\end{document}
